\documentclass[11pt,a4paper]{article}
\usepackage{amsmath,amssymb,amsthm}
\usepackage{geometry}
\usepackage{hyperref}
\geometry{margin=2.5cm}

\newcommand{\h}{h}                   % bandwidth (scale)
\newcommand{\km}{k_m}                % kernel mass
\newcommand{\teff}{T_{\mathrm{eff}}} % effective time
\newcommand{\neff}{n_{\mathrm{eff}}} % effective count
\newcommand{\cps}{\,[\mathrm{cps}]}
\newcommand{\usvh}{\,[\mu\mathrm{Sv/h}]}

\title{Mathematical Description of \texttt{KernelDensityEstimator}}
\author{}
\date{}

\begin{document}
\maketitle
\tableofcontents
\bigskip

\noindent
This document describes the mathematics underlying each method of the
\texttt{KernelDensityEstimator} Kotlin class.  The estimator converts a
stream of Geiger--Müller counts into a continuous dose-rate curve
$\hat\lambda(t)$ in $\mu\mathrm{Sv/h}$, together with Garwood Poisson
confidence intervals.

% ---------------------------------------------------------------------------
\section{Epanechnikov Kernel}

The class uses the \emph{Epanechnikov} kernel with bandwidth $\h > 0$:
\begin{equation}
  K_\h(s) = \frac{3}{4\h}\!\left(1 - \frac{s^2}{\h^2}\right)
             \mathbf{1}_{\{|s|\le\h\}}.
  \label{eq:kernel}
\end{equation}
It is a symmetric, non-negative, compactly supported kernel that integrates to
one and minimises asymptotic mean integrated squared error among all kernels.

\paragraph{Kernel CDF.}
The antiderivative of $K_\h$, needed for interval convolution (Section~\ref{sec:intervals}), is
\begin{equation}
  F_K(x) = \int_{-\infty}^{x} K_\h(s)\,ds =
  \begin{cases}
    0 & x < -\h, \\[4pt]
    \dfrac{3}{4}\!\left(\dfrac{x}{\h} - \dfrac{x^3}{3\h^3} + \dfrac{2}{3}\right)
      & |x| \le \h, \\[4pt]
    1 & x > \h.
  \end{cases}
  \label{eq:Fk}
\end{equation}

\paragraph{Helper \texttt{epanechnikovIntegral}$(u_{\ell}, u_r)$.}
Returns the partial mass of the kernel over a normalised sub-interval:
\begin{equation}
  \texttt{epanechnikovIntegral}(u_\ell, u_r)
    = \int_{u_\ell}^{u_r} \tfrac{3}{4}(1 - u^2)\,du
    = \tfrac{3}{4}\!\left[u - \tfrac{u^3}{3}\right]_{u_\ell}^{u_r}.
  \label{eq:epInt}
\end{equation}
On full support $[-1,1]$ this equals $1$.  It is used for boundary correction
(Section~\ref{sec:boundary}).

\paragraph{Helper \texttt{epanechnikovSquaredIntegral}$(u_{\ell}, u_r)$.}
Returns the squared-kernel integral needed for $\teff$:
\begin{equation}
  \texttt{epanechnikovSquaredIntegral}(u_\ell, u_r)
    = \int_{u_\ell}^{u_r}\!\left(\tfrac{3}{4}\right)^2(1-u^2)^2 du
    = \tfrac{9}{16}\!\left[u - \tfrac{2u^3}{3} + \tfrac{u^5}{5}\right]_{u_\ell}^{u_r}.
  \label{eq:ep2Int}
\end{equation}
On full support $[-1,1]$ this equals $3/5$.

% ---------------------------------------------------------------------------
\section{KDE Estimate}

Let the data consist of $N$ point events at times $t_i$ with counts $n_i$,
and independently of $L$ uniformly distributed sample intervals
$[a_\ell, b_\ell]$ with counts $c_\ell$ (Section~\ref{sec:intervals}).
Let $t_{\mathrm{start}}$ and $t_{\mathrm{end}}$ be the overall data
extents, covering both point events and interval endpoints.

The boundary-corrected KDE estimate at query time $t$ is
\begin{equation}
  \hat\lambda(t)
    = \frac{\displaystyle\sum_{i=1}^{N} n_i\, K_\h(t-t_i)
            \;+\; \Lambda_{\mathrm{iv}}(t)}
           {\km(t)},
  \label{eq:kde}
\end{equation}
where $\Lambda_{\mathrm{iv}}(t)$ is the interval contribution (Section
\ref{sec:intervals}) and $\km(t)$ is the boundary-correction mass (Section
\ref{sec:boundary}).  The estimate is in cps; multiplication by the
calibration factor $\alpha$ $[\mu\mathrm{Sv/h}/\mathrm{cps}]$ gives dose rate.

% ---------------------------------------------------------------------------
\section{Point Events: Sliding-Window Accumulator}
\label{sec:point}

\noindent\textit{Implemented in \texttt{estimateDoseRateHelper}.}

\medskip\noindent
For a query $t$, only events with $|t - t_i| \le \h$ contribute.  Introducing
the numerical-stability offset $\tau = t - t_{\mathrm{start}}$ and
$\tau_i = t_i - t_{\mathrm{start}}$, the point-event numerator of
\eqref{eq:kde} is
\begin{align}
  \sum_{|t-t_i|\le\h} n_i K_\h(t-t_i)
  &= \frac{3}{4\h}\sum_{|t-t_i|\le\h} n_i\!\left(1 - \frac{(\tau-\tau_i)^2}{\h^2}\right)
  \nonumber\\
  &= \frac{3}{4\h}\!\left(S_0 - \frac{\tau^2 S_0 - 2\tau S_1 + S_2}{\h^2}\right),
  \label{eq:pointsum}
\end{align}
where $S_k = \sum_{|t-t_i|\le\h} n_i\,\tau_i^k$ for $k=0,1,2$.

As $t$ moves rightward over the sorted query array, the window $[t-\h, t+\h]$
slides, and $S_0, S_1, S_2$ are updated in $O(1)$ per admitted or evicted
event.  Total complexity is $O(N + Q)$ for $Q$ query points.

% ---------------------------------------------------------------------------
\section{Interval Events: Step Decomposition}
\label{sec:intervals}

\noindent\textit{Implemented in \texttt{addSampleInterval} and
\texttt{estimateDoseRateHelper}.}

\medskip\noindent
An interval $[a, b]$ with $c$ counts models a uniform count-rate density
$w = c/(b-a)$ $[\mathrm{cps}]$.  Its exact contribution to the KDE numerator
is
\begin{equation}
  \Lambda_{[a,b]}(t) = w\int_a^b K_\h(t-\tau)\,d\tau
                     = w\bigl[F_K(t-a) - F_K(t-b)\bigr].
  \label{eq:ivContrib}
\end{equation}
This expresses the interval as two \emph{step events}: a rising step at $a$
with weight $+w$ and a falling step at $b$ with weight $-w$:
\begin{equation}
  \Lambda_{[a,b]}(t) = (+w)\,F_K(t - a) + (-w)\,F_K(t - b).
  \label{eq:stepDecomp}
\end{equation}
For $M$ step events at times $\{t_j\}$ with weights $\{w_j\}$ (so $2L = M$),
the total interval contribution is
\begin{equation}
  \Lambda_{\mathrm{iv}}(t) = \sum_{j=1}^{M} w_j\, F_K(t - t_j).
  \label{eq:ivTotal}
\end{equation}

\subsection{Three-Phase Lifecycle}

As $t$ sweeps rightward, step event $j$ passes through three phases:
\begin{center}
\begin{tabular}{lll}
\hline
Phase & Condition & Contribution \\
\hline
Inactive    & $t < t_j - \h$ & $0$ \\
Transitioning & $|t - t_j| \le \h$ & $w_j F_K(t-t_j)$ \quad (cubic in $t$) \\
Completed   & $t > t_j + \h$ & $w_j \cdot 1$ \quad (constant) \\
\hline
\end{tabular}
\end{center}
The \emph{completed} events contribute only their weights summed into a scalar
$\Sigma_{\mathrm{done}}$, maintained by adding $w_j$ when event $j$ exits the
transitioning window.

\subsection{Transitioning Accumulator Expansion}
\label{sec:tkAccum}

For transitioning events, introduce $\tau_j = t_j - t_{\mathrm{start}}$ and
$\tau = t - t_{\mathrm{start}}$.  Substituting \eqref{eq:Fk} and expanding
$(t - t_j)^3 = (\tau - \tau_j)^3$:
\begin{equation}
  w_j\,F_K(t-t_j)
  = \frac{3}{4}\, w_j \!\left(\frac{\tau - \tau_j}{\h}
      - \frac{(\tau-\tau_j)^3}{3\h^3} + \frac{2}{3}\right).
\end{equation}
Collecting by powers of $\tau_j$ and defining moment accumulators
\begin{equation}
  \tilde T_k = \sum_{\text{transitioning}\,j} w_j\,\tau_j^k,
  \quad k = 0,1,2,3,
  \label{eq:Tk}
\end{equation}
the sum over all transitioning events evaluates to
\begin{equation}
  \Lambda_{\mathrm{trans}}(\tau)
  = \frac{3}{4}\!\left[
      \tilde T_0\!\left(\frac{\tau}{\h} - \frac{\tau^3}{3\h^3} + \frac{2}{3}\right)
    + \tilde T_1\!\left(\frac{\tau^2}{\h^3} - \frac{1}{\h}\right)
    - \tilde T_2\,\frac{\tau}{\h^3}
    + \tilde T_3\,\frac{1}{3\h^3}
  \right].
  \label{eq:transTerm}
\end{equation}
The accumulators $\tilde T_0,\ldots,\tilde T_3$ are updated in $O(1)$ per
admitted or evicted step event, exactly mirroring the $S_0,S_1,S_2$ treatment
of point events.  The total interval contribution is then
\begin{equation}
  \Lambda_{\mathrm{iv}}(t) = \Sigma_{\mathrm{done}} + \Lambda_{\mathrm{trans}}(\tau).
\end{equation}

The sliding window for step events admits event $j$ when
$t_j - \h \le t$ and evicts it (to $\Sigma_{\mathrm{done}}$) when
$t_j + \h \le t$.  Admission is processed before eviction so that a step
event whose entire kernel support has already passed is immediately routed to
$\Sigma_{\mathrm{done}}$ without corrupting the $\tilde T_k$ accumulators.

Overall complexity with $L$ intervals (i.e.\ $M = 2L$ step events) is
$O(N + M + Q)$.

% ---------------------------------------------------------------------------
\section{Boundary Correction}
\label{sec:boundary}

\noindent\textit{Implemented in \texttt{estimateDoseRateHelper} and
\texttt{getConfidenceIntervals}.}

\medskip\noindent
Near the endpoints of the data, the kernel window $[t-\h, t+\h]$ extends
beyond the data range $[t_{\mathrm{start}}, t_{\mathrm{end}}]$.  The
\emph{renormalisation} (boundary) correction divides the KDE numerator by
the fraction of the kernel that lies within the data range.

Using the substitution $u = (t - \tau)/\h$ (so $\tau = t_{\mathrm{start}}$
maps to $u_r = (t - t_{\mathrm{start}})/\h$ and $\tau = t_{\mathrm{end}}$
maps to $u_\ell = (t - t_{\mathrm{end}})/\h$), the boundary-correction kernel
mass is
\begin{equation}
  \km(t) = \int_{t_{\mathrm{start}}}^{t_{\mathrm{end}}} K_\h(t-\tau)\,d\tau
          = \texttt{epanechnikovIntegral}(u_\ell,\, u_r),
  \label{eq:km}
\end{equation}
where
\begin{equation}
  u_\ell = \max\!\left(-1,\;\frac{t - t_{\mathrm{end}}}{\h}\right),\qquad
  u_r    = \min\!\left(+1,\;\frac{t - t_{\mathrm{start}}}{\h}\right).
  \label{eq:ulu}
\end{equation}
For $t$ in the interior (both kernel edges inside the data range),
$u_\ell = -1$, $u_r = +1$, and $\km = 1$.

It is essential that $t_{\mathrm{start}}$ and $t_{\mathrm{end}}$ span
\emph{interval endpoints} as well as point-event timestamps; otherwise the
correction overestimates the boundary effect at the edges of long intervals.

% ---------------------------------------------------------------------------
\section{Confidence Intervals}
\label{sec:ci}

\noindent\textit{Implemented in \texttt{getConfidenceIntervals}.}

\subsection{Effective Observation Time}

Locally, the KDE behaves as if $\neff$ counts were observed over an
equivalent observation window of duration $\teff$.  By matching the
asymptotic Poisson variance of the KDE to that of a direct observation,
\begin{equation}
  \teff = \frac{\km(t)^2}{\displaystyle\int_{t_{\mathrm{start}}}^{t_{\mathrm{end}}}
             K_\h^2(t-\tau)\,d\tau}.
  \label{eq:teff}
\end{equation}
The denominator is computed by substituting $u=(t-\tau)/\h$, giving
$d\tau = -\h\,du$:
\begin{equation}
  \int_{t_{\mathrm{start}}}^{t_{\mathrm{end}}} K_\h^2(t-\tau)\,d\tau
  = \frac{1}{\h}\;\texttt{epanechnikovSquaredIntegral}(u_\ell, u_r),
\end{equation}
so that
\begin{equation}
  \teff = \frac{\h\;\km(t)^2}{\texttt{epanechnikovSquaredIntegral}(u_\ell, u_r)}.
  \label{eq:teffCode}
\end{equation}
For an interior point with full kernel support, this reduces to
$\teff = \h \cdot 1^2 / (3/5) = 5\h/3$.

\subsection{Effective Count and Garwood Interval}

The effective count is
\begin{equation}
  \neff = \hat\lambda(t)\;\teff,
  \label{eq:neff}
\end{equation}
where $\hat\lambda(t)$ is the cps estimate from \eqref{eq:kde}.

The \emph{Garwood} exact Poisson confidence interval at confidence level
$1-\alpha$ for an observed integer count $n$ is
\begin{equation}
  \left(\frac{\chi^2_{2n,\,\alpha/2}}{2},\quad
        \frac{\chi^2_{2(n+1),\,1-\alpha/2}}{2}\right),
\end{equation}
expressed here as bounds on the true count.  Dividing by $\teff$ and
multiplying by the calibration factor $\alpha$ converts to $\mu\mathrm{Sv/h}$.

Because $\neff$ is in general non-integer, the bounds are linearly
interpolated between the two bracketing integers
$n_F = \lfloor\neff\rfloor$ and $n_C = n_F + 1$ using weight
$\xi = \neff - n_F$:
\begin{align}
  \chi_L &= \chi^2_{2n_F,\,0.025}
           + \xi\!\left(\chi^2_{2n_C,\,0.025}
                        - \chi^2_{2n_F,\,0.025}\right), \\[4pt]
  \chi_R &= \chi^2_{2(n_F+1),\,0.975}
           + \xi\!\left(\chi^2_{2(n_C+1),\,0.975}
                        - \chi^2_{2(n_F+1),\,0.975}\right).
  \label{eq:chiInterp}
\end{align}
The dose-rate confidence bounds are
\begin{equation}
  \lambda^{\pm}\usvh
  = \frac{\chi_{L,R}}{2\,\teff}\cdot\alpha.
  \label{eq:ciBounds}
\end{equation}

% ---------------------------------------------------------------------------
\section{Summary of Complexity}

\begin{center}
\begin{tabular}{lll}
\hline
Operation & Per-call cost & Amortised total \\
\hline
\texttt{addPoint} & $O(1)$ & $O(N)$ \\
\texttt{addSampleInterval} & $O(1)$ & $O(L)$ \\
\texttt{estimateDoseRateHelper} & $O(N + M + Q)$ & --- \\
\texttt{getConfidenceIntervals} & $O(N + M + Q)$ & --- \\
\hline
\end{tabular}
\end{center}
$N$ = number of point events, $L$ = number of intervals,
$M = 2L$ = number of step events, $Q$ = number of query points.

\end{document}
